% SPDX-License-Identifier: MIT
\chapter{Editing the tree}
\label{ch:editing}

Every operation in this chapter is undoable (\keys{Ctrl+Z}) and every one is also
available from the Python API (Chapter~\ref{ch:python}).

\section{Rooting}

Where the root sits changes what the figure claims. It does not change the tree's
topology in the unrooted sense, but it entirely changes which groups look
monophyletic, so it is worth being deliberate about.

\begin{center}
\begin{tabular}{@{}llp{0.48\linewidth}@{}}
\toprule
\textbf{Operation} & \textbf{Key} & \textbf{Description} \\
\midrule
Reroot here      & \keys{R} & Place the root on the branch above the selected node. \\
Midpoint root    & \keys{M} & Root at the midpoint of the longest tip-to-tip path. \\
Root on outgroup & \keys{Ctrl+Shift+R} & Root against a named taxon or set of taxa. \\
Unroot           & --- & Remove the root, collapsing it into a trifurcation. \\
\bottomrule
\end{tabular}
\end{center}

\textbf{Midpoint rooting} assumes a roughly constant rate of evolution across the
tree. It is a reasonable default when you have no outgroup, and a poor one when
rates vary substantially between lineages.

\textbf{Outgroup rooting} is preferable whenever you have a defensible outgroup,
because it encodes an actual biological hypothesis rather than an assumption
about rates.

\galleryfig{feature-midpoint-rooted}{The example tree after midpoint rooting.}{1.0}

\begin{quote}
\textbf{Branch lengths survive rerooting.} A branch length belongs to the
\emph{edge} between a node and its parent, and rerooting reverses the direction
of some edges. \app{} carries the lengths across those reversals, and the test
suite asserts that all pairwise tip-to-tip distances are unchanged by a reroot.
Rerooting is a change of viewpoint, not of data.
\end{quote}

\section{Ordering}

Reordering changes only the cross coordinate, so it can never alter a distance in
the figure.

\begin{center}
\begin{tabular}{@{}llp{0.48\linewidth}@{}}
\toprule
\textbf{Operation} & \textbf{Key} & \textbf{Description} \\
\midrule
Ladderize ascending  & \keys{L}       & Smaller subtrees first, giving the characteristic staircase. \\
Ladderize descending & \keys{Shift+L} & Larger subtrees first. \\
Rotate children      & \keys{Ctrl+R}  & Swap the order of the selected node's children. \\
\bottomrule
\end{tabular}
\end{center}

Ladderizing is almost always worth doing before exporting: it makes the branching
order far easier to follow, and costs nothing in accuracy.

\section{Collapsing}

Collapsing hides a clade's internal structure while keeping its position and
extent, which is the standard way to show a large tree without printing every
tip.

\keys{C} collapses or expands the selected node; \keys{Shift+C} expands
everything. A collapsed clade is drawn as a triangle whose depth reflects the
subtree it stands for.

\galleryfig{feature-collapsed-clade}{A collapsed clade. The triangle's size is
derived from the hidden subtree rather than being fixed, so the figure still
conveys how much was folded away.}{1.0}

The number of rows a collapsed clade occupies is configurable:
\opt{collapse\_rows\_mode} may be a square-root function of the subtree size (the
default, which keeps large clades from dominating), a fixed number of rows, or
proportional. \opt{collapse\_shape} selects the glyph.

\section{Pruning and extracting}

\begin{center}
\begin{tabular}{@{}llp{0.48\linewidth}@{}}
\toprule
\textbf{Operation} & \textbf{Key} & \textbf{Description} \\
\midrule
Remove selected taxa & \keys{Del}          & Delete tips and tidy away the internal nodes left dangling. \\
Extract subtree      & \keys{Ctrl+Shift+X} & Replace the document with the selected clade alone. \\
\bottomrule
\end{tabular}
\end{center}

When pruning removes a tip, the internal node above it may be left with a single
child. Such nodes are collapsed away, and the branch lengths on either side are
\emph{summed} rather than discarded, so the distances between the remaining tips
are preserved.

\section{Selecting}

Editing acts on the selection, so the selection tools matter.

\begin{itemize}
  \item Click a node, or drag a rubber band across several.
  \item \keys{Ctrl+Shift+A} extends the selection to the entire clade below the
        current node.
  \item The Search panel (\keys{Ctrl+F}) finds taxa by name or regular
        expression and can select every match at once --- the practical way to
        select a scattered group.
  \item \keys{Esc} clears the selection.
\end{itemize}
